\section{o1 的启发：RL、CoT 与反思 pattern}

上一章指出 next token prediction 的缺陷，本章看张祥雨如何理解 o1。访谈中他的判断很明确：o1 的关键不只是用了 RL，而是激发了一种思维链 pattern。传统 RL 的许多算法名字并不是重点；重点是，模型在预训练中已经见过许多思考、反思、检查和重做的模式，RL 可以把这些模式强化出来。

\begin{figure}[H]
\centering
\includegraphics[width=0.96\textwidth]{rl-cot-pattern.png}
\caption{RL 与 CoT Pattern：RL 不只是算法名，关键是激发可搜索的思考模式。自制概念图，依据 00:50:48--01:01:52 对谈内容整理。}
\end{figure}

\begin{knowledgebox}{读图：RL 不是魔法，pattern 才是杠杆}
图中 Pretrain 提供已有模式，RL 用目标反馈强化，CoT 把步骤展开，Reflection 修正路径，Critical Decision 处理关键分支。o1 的启发在于：模型不是凭空学会思考，而是把预训练里稀疏存在的思考 pattern 放大。
\end{knowledgebox}

\subsection{Critical Decision：一个 token 解决不了的地方}

本节解释为什么反思重要。推理过程中不是每个 token 都同样关键。很多 token 只是表达格式，真正决定成败的是少数 critical decision：走左边还是右边、用哪个公式、是否需要回头检查、某个中间结论是否可靠。当一个分支选择的复杂度超过单个 token 的表达能力时，模型需要展开更多步骤，把决策问题拆成可搜索空间。

\begin{figure}[H]
\centering
\includegraphics[width=0.96\textwidth]{meta-cot-critical-decision.png}
\caption{Meta-CoT 与 Critical Decision：当一个 token 解决不了分支选择，就需要反思和元思维链。自制概念图，依据 01:01:52--01:10:57 对谈内容整理。}
\end{figure}

\begin{knowledgebox}{读图：Meta-CoT 是 CoT 的 CoT}
从 State 到 Branch，再到 One token?、Reflect 和 Meta-CoT，图中表达的是“对思考过程再思考”。模型不只是写步骤，而是在关键节点判断当前路径是否足够、是否要回退、是否要换分支。
\end{knowledgebox}

\begin{importantbox}{课堂提示：动作空间比算法名字更关键}
张祥雨反复强调，很多时候不是具体 RL 算法多神秘，而是动作空间是否包含正确动作。没有反思动作，模型很难在关键决策点搜索；有了反思动作，RL 才有东西可强化。
\end{importantbox}

\subsection{本章小结}

o1 给多模态的启发是：智能不只是模型更大，还要有可搜索的中间过程。CoT、反思和 Meta-CoT 把不可解的单 token 决策，变成可以分解、试错和优化的动作空间。

